\chapter{Sheaves}

\section*{PRESHEAVES}

\begin{definition}
	Let $X$ be a topological space. A \textcolor{orange}{presheaf} of sets is a functor $\mathcal{F}:\mathbf{Open}(X)^\text{op}\rightarrow\mathbf{Set}$, where $\operatorname{ob}\left(\mathbf{Open}(X)\right)$ is the set of all open subsets of $X$, and for any $U, V \in \operatorname{ob}\left(\mathbf{Open}(X)\right)$,
\begin{equation*}
	\operatorname{Hom}\left(U,V\right)=\left\{\begin{aligned}
		&\left\{\iota_{VU}\right\}, && \text{if }V\subseteq U;\\
		&\varnothing, && \text{if }V\nsubseteq U.
	\end{aligned}\right.
\end{equation*}
where $\iota_{UV}$ denotes the unique inclusion morphism from $V$ to $U$.

The category $\mathbf{Set}$ in the definition of a presheaf can be replaced by an arbitrary category $\mathcal{C}$, and the corresponding functor is called a \textcolor{orange}{$\mathcal{C}$-valued presheaf}. 
\end{definition}

\begin{definition}
	An element $s \in \mathcal{F}(U)$ is called a \textcolor{orange}{section} of $\mathcal{F}$ over $U$.
	
	The set $\mathcal{F}(U)$ is called the \textcolor{orange}{set of sections} (or simply the \textcolor{orange}{sections}) over $U$.
	
	The map $\mathcal{F}\!\left(\iota_{VU}\right)$ is called the \textcolor{orange}{restriction map}, and $\mathcal{F}\!\left(\iota_{VU}\right)(s)$ is written to be $s|_V$ if $s\in\mathcal{F}(U)$.
\end{definition}

\begin{definition}
	Let $\mathcal{F}$ and $\mathcal{G}$ be presheaves on $X$. A \textcolor{orange}{morphism} $\varphi:\mathcal{F}\rightarrow\mathcal{G}$ is a collection of maps $\varphi_U:\mathcal{F}(U)\rightarrow\mathcal{G}(U)$ for each open set $U \subseteq X$, such that whenever $V\subseteq U$, the diagram
	\begin{equation*}
		\begin{tikzcd}
			\mathcal{F}(U)\arrow[r,"\mathcal{F}\left(\iota_{VU}\right)"]\arrow[d,"\varphi_U"]&[2em]\mathcal{F}(V)\arrow[d,"\varphi_V"]\\
			\mathcal{G}(U)\arrow[r,"\mathcal{G}\left(\iota_{VU}\right)"]&\mathcal{G}(V)
		\end{tikzcd}
	\end{equation*}
	commutes. All presheaves on $X$ form a category, denoted by \textcolor{orange}{$\mathbf{PSh}(X)$}.
	
	An isomorphism is defined to be a morphism having a two-sided inverse.
\end{definition}

\begin{definition}
	Let $\varphi:\mathcal{F}\rightarrow\mathcal{G}$ be a morphism of presheaves. The \textcolor{orange}{presheaf kernel} of $\varphi$, \textcolor{orange}{presheaf cokernel} of $\varphi$, and \textcolor{orange}{presheaf image} of $\varphi$ are defined to be the presheaves given by $U\mapsto\ker(\varphi_U)$, $U\mapsto\operatorname{coker}(\varphi_U)$ and $U\mapsto\operatorname{im}(\varphi_U)$, respectively.
\end{definition}

\begin{definition}
	Let $\mathcal{F}$ be a presheaf on $X$. The \textcolor{orange}{stalk} of $\mathcal{F}$ at $x \in X$ is the direct limit
	\begin{equation*}
		\mathcal{F}_x=\varinjlim_{U\ni x}\mathcal{F}(U).
	\end{equation*}
	
	The elements of a stalk $\mathcal{F}_x$ are the \textcolor{orange}{germs} of sections of $\mathcal{F}$ at point $x$.
	
	Denote $\mathcolor{orange}{s_x}\in\mathcal{F}_x$ to be the image of $s\in\mathcal{F}(U)$ under the canonical morphism of the direct limit.
\end{definition}

\begin{definition}
	The \textcolor{orange}{\'etal\'e space} of the presheaf $\mathcal{F}$ is the disjoint union of all stalks
	\begin{equation*}
		E(\mathcal{F})=\coprod_{x\in X}\mathcal{F}_x
	\end{equation*}
	equipped with the topology generated by the basis of sets $U_s=\left\{s_x\in\mathcal{F}_x\mid x\in U\right\}$ for all open sets $U\subseteq X$ and all sections $s\in\mathcal{F}(U)$, where $s_x$ denotes the image of $s$ in the stalk $\mathcal{F}_x$.
	
	There is a canonical projection $\pi: E(\mathcal{F}) \rightarrow X$ defined by $s_x \mapsto x$ for each $s_x \in \mathcal{F}_x$. The topology of the \'etal\'e space ensures that $\pi$ is a continuous map (more precisely, a local homeomorphism).
\end{definition}


\section*{SHEAVES}

\begin{definition}
	A presheaf $\mathcal{F}$ on a topological space $X$ is a \textcolor{orange}{sheaf} if, for any open set $U\subseteq X$ and any open covering $\left\{V_i\right\}_{i\in I}$ of $U$, the following conditions hold:
	\begin{enumerate}[label=\textup{(\arabic*)}]
		\item \textcolor{red}{Locality:} For any sections $s, t\in\mathcal{F}(U)$, if $s|_{V_i}=t|_{V_i}$ for all $i\in I$, then $s = t$.
		\item \textcolor{red}{Gluing:} For a family of sections $\left\{s_i\in\mathcal{F}\!\left(V_i\right)\right\}_{i\in I}$ satisfying $s_i|_{V_i\cap V_j}=s_j|_{V_i\cap V_j}$ for all $i, j\in I$, there exists a section $s\in\mathcal{F}(U)$ such that $s|_{V_i}=s_i$ for all $i\in I$.
	\end{enumerate}
	
	The category of sheaves on $X$, denoted by \textcolor{orange}{$\mathbf{Sh}(X)$}, is defined as the full subcategory of $\mathbf{PSh}(X)$ consisting of those presheaves that satisfy both the locality and gluing properties.
\end{definition}

\begin{proposition}{1.1}
	Let $\varphi:\mathcal{F}\rightarrow\mathcal{G}$ be a morphism of presheaves on a topological space $X$. If $\varphi$ is an isomorphism, then the induced map $\varphi_x:\mathcal{F}_x\rightarrow\mathcal{G}_x$ is an isomorphism for every $x\in X$. If in addition $\mathcal{F}$ and $\mathcal{G}$ are sheaves, then $\varphi$ is an isomorphism iff $\varphi_x$ is an isomorphism for all $x\in X$.
\end{proposition}

\begin{definition}
	Let $\mathcal{F}$ be a presheaf. The \textcolor{orange}{associated sheaf} (or \textcolor{orange}{sheafification}) of $\mathcal{F}$ is a sheaf $\mathcal{F}^+$ which maps each open set $U$ to the set of maps $s:U\rightarrow E(\mathcal{F})$, such that
	\begin{enumerate}[label=\textup{(\arabic*)}]
		\item for each $x\in U$, $s(x)\in\mathcal{F}_x$;
		\item for each $x\in U$, there exists an open neighborhood $V\subseteq U$ of $x$ and a section $t\in\mathcal{F}(V)$ such that $t_y=s(y)$ for all $y\in V$.
	\end{enumerate}
	
	There is a canonical morphism of presheaves $\theta:\mathcal{F}\rightarrow\mathcal{F}^+$ defined by $\mathcal{F}(U)\mapsto\mathcal{F}^+(U)$.
\end{definition}

\begin{proposition}{1.2}[Universal property of sheafification]
	Let $\theta:\mathcal{F}\rightarrow\mathcal{F}^+$ be the canonical morphism. Then for any sheaf $\mathcal{G}$ and any morphism $\varphi:\mathcal{F} \rightarrow\mathcal{G}$, there exists a unique morphism of presheaves $\psi:\mathcal{F}^+\rightarrow\mathcal{G}$ such that $\varphi=\psi\circ\theta$.
\end{proposition}

\begin{definition}
	Let $\varphi:\mathcal{F}\rightarrow\mathcal{G}$ be a morphism of sheaves. The \textcolor{orange}{kernel} of $\varphi$ is defined to be $\ker_\mathbf{PSh}\varphi$; the \textcolor{orange}{cokernel} and \textcolor{orange}{image} of $\varphi$ are defined to be $(\operatorname{coker}_\mathbf{PSh}\varphi)^+$ and $(\operatorname{im}_\mathbf{PSh}\varphi)^+$, respectively.
	
	The morphism $\varphi$ is \textcolor{orange}{injective} if $\ker\varphi=0$. It is \textcolor{orange}{surjective} if $\operatorname{im}\varphi=\mathcal{G}$.
	
	Let $\mathcal{F}^\prime$ be a subsheaf of $\mathcal{F}$. The \textcolor{orange}{quotient sheaf} $\mathcal{F}/\mathcal{F}^\prime$ is defined to be the associated sheaf of the presheaf $U\mapsto\mathcal{F}(U)/\mathcal{F}^\prime(U)$. Use the definition of associated sheaf, it follows that $(\mathcal{F}/\mathcal{F}^\prime)_x=\mathcal{F}_x/\mathcal{F}_x^\prime$.
\end{definition}

\begin{definition}
	Let $f:X\rightarrow Y$ be a continuous map of topological spaces.
	
	The \textcolor{orange}{direct image} of a sheaf $\mathcal{F}$ on $X$ is the sheaf $f_*\mathcal{F}$ on $Y$ defined by $f_*\mathcal{F}(V)=\mathcal{F}\!\left(f^{-1}(V)\right)$.
	
	The \textcolor{orange}{inverse image} of a sheaf $\mathcal{G}$ on $Y$ is the sheaf $f^{-1}\mathcal{G}$ on $X$ defined by
	\begin{equation*}
		f^{-1}\mathcal{G}(U)=\left(\varinjlim_{V\supseteq f(U)}\mathcal{G}(V)\right)^+.
	\end{equation*}
	
	The \'etal\'e space of $f^{-1}\mathcal{G}$ is the pullback $E(\mathcal{G})\times_YX$ satisfying the commutative diagram
	\begin{equation*}
		\begin{tikzcd}
			E(\mathcal{G})\times_YX\arrow[r]\arrow[d]&E(\mathcal{G})\arrow[d,"\pi_Y"]\\
			X\arrow[r,"f"]&Y
		\end{tikzcd}
	\end{equation*}
\end{definition}

\begin{definition}
	Let $\mathcal{F}$ be a sheaf on $X$ and $Z$ is a subset of $X$. The \textcolor{orange}{restriction} of $\mathcal{F}$ to $Z$ is the inverse image $i^{-1}\mathcal{F}$, where $i:Z\rightarrow X$ is the inclusion map, denote by \textcolor{orange}{$\mathcal{F}|_Z$}.
\end{definition}